\documentclass[11pt]{article}
\usepackage[utf8]{inputenc}
\usepackage[T1]{fontenc}
\usepackage[italian]{babel}
\usepackage{geometry}
\geometry{a4paper, margin=1in}
\usepackage{xcolor}
\usepackage{listings}
\usepackage{hyperref}

\usepackage{tcolorbox}
\tcbuselibrary{breakable}
\begin{document}

\begin{tcolorbox}[colback=blue!5!white,colframe=blue!75!black,title=\textbf{DOMANDA}, breakable]
Esercizio 3\\
Modificare il sorgente del codice in modo da cambiare nome ai seguenti eventi:
\begin{itemize}
    \item \texttt{message} $\rightarrow$ \texttt{messaggio}
    \item \texttt{UploadChat} $\rightarrow$ \texttt{aggiornamento}
\end{itemize}
\end{tcolorbox}

\begin{tcolorbox}[colback=green!5!white,colframe=green!75!black,title=\textbf{RISPOSTA}, breakable]
\textbf{Esercizio 36 - Rinomina Eventi WebSocket}

\textbf{Modifiche effettuate al codice sorgente:}

\begin{enumerate}
    \item \textbf{Lato Server (\texttt{server.js}):}
    \begin{itemize}
        \item Modificato l'evento in ascolto da \texttt{'message'} a \texttt{'messaggio'}:\\
        \texttt{webSocket.on('messaggio', function(data) \{ ... \});}
        \item Modificato l'evento emesso in broadcast da \texttt{'UploadChat'} a \texttt{'aggiornamento'}:\\
        \texttt{io.sockets.emit('aggiornamento', data);}
    \end{itemize}

    \item \textbf{Lato Client (\texttt{public/chat.js}):}
    \begin{itemize}
        \item Modificato il nome dell'evento inviato al click del pulsante da \texttt{'message'} a \texttt{'messaggio'}:\\
        \texttt{webSocket.emit('messaggio', \{ ... \});}
        \item Modificato il nome dell'evento in ascolto per ricevere i messaggi da \texttt{'UploadChat'} a \texttt{'aggiornamento'}:\\
        \texttt{webSocket.on('aggiornamento', function(data) \{ ... \});}
    \end{itemize}
\end{enumerate}
\end{tcolorbox}

\end{document}
